% Generated draft: requires mathematical and source review.
\documentclass[11pt]{article}
\usepackage[margin=25mm]{geometry}
\usepackage{fontspec}
\setmainfont{Latin Modern Roman}
\usepackage{amsmath,amssymb,mathtools,unicode-math}
\setmathfont{Latin Modern Math}
\usepackage{enumitem,needspace,xurl,hyperref}
\setlength{\parindent}{0pt}
\setlength{\parskip}{5pt}
\setlength{\emergencystretch}{3em}
\newcommand{\N}{\mathbb{N}}
\newcommand{\Z}{\mathbb{Z}}
\newcommand{\Q}{\mathbb{Q}}
\newcommand{\R}{\mathbb{R}}
\newcommand{\C}{\mathbb{C}}
\newcommand{\F}{\mathbb{F}}
\newcommand{\A}{\mathbb{A}}
\newcommand{\PP}{\mathbb{P}}
\newcommand{\E}{\mathbb{E}}
\newcommand{\eps}{\varepsilon}
\newcommand{\dd}{\,\mathrm{d}}
\begin{document}
\section*{UnsolvedMath ID: 30004583}
\textbf{Source problem code:} catalogue entry
\subsection{Definitions and mathematical statement}
The supplied record concerns arithmetic education, textbook content, systems of weights, measures and currency, and resources for everyday numerical computation during the ``long nineteenth century.'' It does not specify the dates denoted by this period, the geographical scope, the populations to be studied, or criteria for measuring educational availability and computational skills.

\textbf{Research question.} In the long nineteenth century, to what extent was arithmetic education available and to whom; what forms did it take and what skills did textbooks impart; what systems of weights, measures, and currency existed; and what resources supported everyday numerical computation?

This is a historiographical research program, not a supplied mathematical proposition. No formal domains, quantitative measures, hypotheses, boundary conditions, or mathematical quantifiers are provided. A precise mathematical assertion cannot be extracted without adding assumptions absent from the source. The research question is therefore retained with these qualifications unresolved.
\subsection{Short English statement}
Investigate who could obtain arithmetic education during the ``long nineteenth century,'' its forms and textbook skills, the systems of weights, measures and currency in use, and the resources supporting everyday calculation. The dates, geographical coverage and assessment criteria remain unspecified. The supplied record labels the question open as a historiographical program; this draft does not independently verify that assessment.
\subsection{Sources}
\begin{itemize}
\item Local registry: supplied UnsolvedMath record, permanent ID \(30004583\), title \emph{Arithmetic Education and Numerical Resources in the Nineteenth Century}; source problem code \texttt{catalogue entry}; collection: Oberwolfach Reports Open Problems. No individual proposer is identified. The recorded proposed year is 2020.
\item Supplied report citation: \emph{History of Mathematics: A Global Cultural Approach} (2021), Oberwolfach Reports, pp. 2021--2092, DOI \texttt{10.4171/owr/2020/41}. \url{https://doi.org/10.4171/owr/2020/41}. The background attributes the collection to Oberwolfach Reports, which document mathematical workshops and their open-problem discussions.
\item Supplied literature source, labelled ``original Oberwolfach report'': \url{https://doi.org/10.4171/owr/2020/41}.
\item Literature assessment dated exactly \texttt{2026-08-22}: the supplied curation labels the problem ``open'' and describes it as an open historiographical program encompassing historical education, textbooks, measures, currency and computational resources, rather than a settled claim. This is supplied dated curation, not independently verified literature assessment. The background explicitly calls for independent verification of time-sensitive status and cited claims.
\item The registry marks the statement ``exact'' and reports that it is internally readable and complete and was retained byte-for-byte. These are supplied verification claims; the report and its wording have not been independently checked for this draft.
\end{itemize}
\end{document}
